\subsection{Theme Synthesis — modelの外側にworkflow-awareな層が生まれる}

4月のserving群は、後の「agent runtime」という大きな言葉を先取りするより、何が従来と違うのかを具体的に見る方が有益だ。複数callの依存関係、context-aware resource control、workflow semantics、KV restoration、multimodal MoE routing、そしてmodel supportの速い更新。これらはmodel parameterそのものではなく、実行時に初めて現れる条件である。5月以降にruntime全体の競争が強まる前に、4月にはserving layerがapplication/workflowの性質を意識し始めていた。\autocite{src-0fc9d11483aa,src-2a6eae37bd39,src-551442a9e215,src-b790ec92aa01,src-d91793d0a6fa,src-e1e959857ced,src-fe12f92c05d7}


\begin{claimboundary}
Claim Boundary: Scepsy、KAIROS、CacheFlowなどの比較・性能結果はpaper authorsによる報告であり、本稿では独立再現していない。設計上の問題設定を一般化しても、報告されたgainを任意環境の性能保証へ変換しない。\autocite{src-0fc9d11483aa,src-551442a9e215,src-e1e959857ced}
\end{claimboundary}

\begin{claimboundary}
Claim Boundary: vLLM release notesはrelease chronologyとprojectが記録したimplementation／model support変更を確認する資料である。performanceとcompatibilityの一般保証ではない。\autocite{src-b790ec92aa01,src-fe12f92c05d7}
\end{claimboundary}

\begin{claimboundary}
Chronology Boundary: later revisionを保持する論文ではApril初回投稿と後日改訂を分離する。ReaLBは2026年5月9日更新、Pythiaは2026年5月14日更新のlocatorを保持しているため、post-Aprilの本文変更を4月の状態として引用しない。\autocite{src-2a6eae37bd39,src-d91793d0a6fa}
\end{claimboundary}
